\renewcommand{\thetheorem}{17.c.\arabic{theorem}}
\setcounter{theorem}{0}

% ---------------------------------------------------------------------------
% Provenance of the prose in this file.
%   % Extractor: <model>   the block below was transcribed from Prof. Biran's
%                          handwritten notes by that model
%   % Creator:   <model>   the block below was written by that model and has no
%                          counterpart in the notes
% A tag holds until the next tag.
% ---------------------------------------------------------------------------

% Extractor: Gemini 3.1 Pro
\section{Column Rank and Row Rank}
\label{sec:column_row_rank}

\begin{theorem}[Equality of Column and Row Rank]
% thm:17.c.1
\label{thm:rank_row_column_equality}
Let $A \in \M_{m \times n}(K)$. Then, $\rank_{\text{col}}(A) = \rank_{\text{row}}(A)$.
\end{theorem}

\begin{remark}[Notation and Properties of Rank]
\begin{enumerate}
    \item Because of this theorem, we will from now on denote $\rank(A) := \rank_{\text{col}}(A) = \rank_{\text{row}}(A)$.
    \item Suppose $A'$ is a row-reduced echelon matrix that is row-equivalent to $A$. Then, we have:
    \begin{equation*}
        \rank(A) = \rank_{\text{row}}(A) = \rank_{\text{row}}(A') = \text{number of pivots in } A'
    \end{equation*}
    \item Recall that for a linear map $T: V \to W$, we defined $\rank(T) := \dim \Img(T)$. Moreover, if $A \in \M_{m \times n}(K)$, then $\rank(T_A) = \dim \Img(T_A) = \rank_{\text{col}}(A)$.
\end{enumerate}
\end{remark}

\subsection{Preparation for the Proof}

\begin{lemma}[Rank Invariance under Composition with Isomorphisms]
% lemma:17.c.2
\label{lemma:rank_invariance_isomorphisms}
Let $T: V \to W$ be a linear map, and let $S: W \to W'$ and $L: P \to V$ be isomorphisms. Then:
\begin{enumerate}[label=\textbf{(\alph*)}]
    \item $\rank(S \circ T) = \rank(T)$.
    \item $\rank(T \circ L) = \rank(T)$.
\end{enumerate}
\end{lemma}

\begin{proof}
\begin{enumerate}[label=\textbf{(\alph*)}]
    \item By definition, $\rank(S \circ T) = \dim (S(T(V)))$. Since $S$ is an isomorphism, its restriction to the subspace $T(V) \subseteq W$ is injective. An injective linear map preserves the dimension of a vector space, and therefore:
    \begin{equation*}
        \dim (S(T(V))) = \dim (T(V)) = \rank(T)
    \end{equation*}

    \item By definition, $\rank(T \circ L) = \dim (T(L(P)))$. Since $L$ is an isomorphism, it is surjective, which implies that $L(P) = V$. Consequently:
    \begin{equation*}
        \dim (T(L(P))) = \dim (T(V)) = \rank(T)
    \end{equation*}
\end{enumerate}
\end{proof}

% Extractor: Opus 5 (High)
\begin{lemma}[Invertibility and Isomorphism]
% lemma:17.c.3
\label{lemma:invertible_iff_T_A_iso}
Let $A \in \M_{n \times n}(K)$. Then $A$ is invertible if and only if $T_A: K^n \to K^n$ is an isomorphism.
\end{lemma}

\begin{proof}
Suppose first that $A$ is invertible. Then $T_{A^{-1}} = (T_A)^{-1}$, so $T_A$ is an isomorphism.

Conversely, suppose that $T_A$ is an isomorphism. We know that $A = [T_A]_{\mathcal{E}_n}^{\mathcal{E}_n}$, and by a previous result the representation matrix of an isomorphism is invertible, so $A$ is invertible.
\end{proof}

% Creator: Opus 5 (High)
\begin{remark}[Invertibility of the Transpose]
\label{rem:transpose_invertibility}
It is worth recording alongside this that $A \in \GL_n(K)$ if and only if $A\transp \in \GL_n(K)$, with $(A\transp)^{-1} = (A^{-1})\transp$. This is \cref{exc:transpose_properties_multiplication} \textbf{(b)} of \cref{sec:matrix_multiplication}, and it needs nothing from the present section: applying $(XY)\transp = Y\transp X\transp$ to $A^{-1}A = A A^{-1} = I_n$ exhibits $(A^{-1})\transp$ as a two-sided inverse of $A\transp$, and the converse follows by applying the same argument to $A\transp$ and using $(A\transp)\transp = A$.
\end{remark}

% Extractor: Gemini 3.1 Pro

\begin{lemma}[Rank Invariance under Matrix Equivalence]
% lemma:17.c.4
\label{lemma:rank_invariance_equivalence}
Let $A, B \in \M_{m \times n}(K)$ such that $B = P A Q$, where $P \in \GL_m(K)$ and $Q \in \GL_n(K)$. Then:
\begin{enumerate}[label=\textbf{(\alph*)}]
    \item $\rank_{\text{col}}(A) = \rank_{\text{col}}(B)$.
    \item $\rank_{\text{row}}(A) = \rank_{\text{row}}(B)$.
\end{enumerate}
\end{lemma}

\begin{proof}
\begin{enumerate}[label=\textbf{(\alph*)}]
    \item Transitioning to the associated linear maps, we have $T_B = T_P \circ T_A \circ T_Q$. Since $P$ and $Q$ are invertible matrices, both $T_P$ and $T_Q$ are isomorphisms. Applying \cref{lemma:rank_invariance_isomorphisms}, we deduce that $\rank(T_B) = \rank(T_A)$. Recalling that the rank of a linear map given by a matrix matches the column rank of that matrix, we obtain $\rank_{\text{col}}(B) = \rank_{\text{col}}(A)$.

    \item Taking the transpose of the equation $B = PAQ$ yields $B\transp = Q\transp A\transp P\transp$. Since $P$ and $Q$ are invertible matrices, their transposes $P\transp$ and $Q\transp$ are also invertible (we have previously seen that $(P\transp)^{-1} = (P^{-1})\transp$). We can now apply statement \textbf{(a)} of this \Cref{lemma:rank_invariance_equivalence}, which we have already proven, to the matrices $A\transp$ and $B\transp$. This gives us $\rank_{\text{col}}(B\transp) = \rank_{\text{col}}(A\transp)$. By the definition of the transpose, this immediately translates to $\rank_{\text{row}}(B) = \rank_{\text{row}}(A)$.
\end{enumerate}
\end{proof}

\subsection{Proof of the Main Theorem}

We are now ready for the proof of \cref{thm:rank_row_column_equality}.

\begin{proof}[Proof of \Cref{thm:rank_row_column_equality}]
Let $r := \rank_{\text{col}}(A)$. By \cref{cor:matrix_equivalence_normal_form}, we know that there exist invertible matrices $P \in \GL_m(K)$ and $Q \in \GL_n(K)$ such that:
\begin{equation*}
    P A Q = \begin{pmatrix} I_r & O \\ O & O \end{pmatrix} =: B
\end{equation*}
By \Cref{lemma:rank_invariance_equivalence}, we know that $\rank_{\text{row}}(A) = \rank_{\text{row}}(B)$. Thus, it suffices to determine the row rank of $B$.

The matrix $B$ is of the block form $\left(\begin{smallmatrix} I_r & O \\ O & O \end{smallmatrix}\right)$. The non-zero rows of $B$ are precisely the first $r$ standard basis vectors $e_1\transp, e_2\transp, \dots, e_r\transp \in K_{\text{row}}^n$. Since these vectors form a linearly independent set, the dimension of the row space of $B$ is exactly $r$. Consequently, $\rank_{\text{row}}(B) = r$.

Therefore, we conclude that $\rank_{\text{row}}(A) = r = \rank_{\text{col}}(A)$, completing the proof.
\end{proof}

\begin{importantremark}[Column Space versus Row Space]
Although $\rank_{\text{col}}(A) = \rank_{\text{row}}(A)$, the column space and the row space are, in general, different vector spaces. First of all, if $A \in \M_{m \times n}(K)$, then the column space is a subspace of $K^m$, whereas the row space is a subspace of $K_{\text{row}}^n$ (which is isomorphic to $K^n$). But even in the case where $m = n$ (i.e., $A \in \M_{n \times n}(K)$), these two spaces are generally distinct! They share the exact same dimension, but they are not the same space.
\end{importantremark}

% Creator: Opus 5 (High)
\begin{ainote}
The transcript has all four numbered statements of this section, but 17.c.3 is
not the one in the notes. Prof. Biran's Lemma 2 records that $A$ is invertible
if and only if $T_A$ is an isomorphism, in two lines; the transcript put in its
place a lemma on the invertibility of $A\transp$, with a proof that runs
through row-reduced echelon forms and pivots. That lemma is true, and its first
sentence is precisely Prof. Biran's Lemma 2 used as a step, but the statement is
not his and nothing in this section uses it. The notes' lemma is restored as
\cref{lemma:invertible_iff_T_A_iso}, and the transpose statement is kept as
\cref{rem:transpose_invertibility}, where it belongs: it is
\cref{exc:transpose_properties_multiplication} \textbf{(b)} of
\cref{sec:matrix_multiplication}, needs nothing from this section, and comes
with the inverse written out.

Two citations were vague, \qt{a previous proposition regarding the
equivalence of matrices} is \cref{cor:matrix_equivalence_normal_form}, and
$\GL$ was typeset as plain $GL$ in the main proof.
\end{ainote}